\chapter*{References}
\addcontentsline{toc}{chapter}{References}
\small
\begin{itemize}[leftmargin=*,label={}]
\item Abbott JK, Morrow EH (2011) Obtaining snapshots of genetic variation using hemiclonal analysis. \emph{Trends Ecol Evol} 26:359--368.
\item Chippindale AK, Gibson JR, Rice WR (2001) Negative genetic correlation for adult fitness between sexes reveals ontogenetic conflict in \emph{Drosophila}. \emph{PNAS} 98:1671--1675.
\item Collet JM \emph{et al.}\ (2016) \emph{Evolution} (PMC5069644); values from the primary text.
\item Comeron JM, Ratnappan R, Bailin S (2012) The many landscapes of recombination in \emph{Drosophila melanogaster}. \emph{PLoS Genet} 8:e1002905.
\item Connallon T, Clark AG (2010) Sex linkage, sex-specific selection, and the role of recombination in the evolution of sexually dimorphic gene expression. \emph{Evolution} 64:3417--3442.
\item Fry JD (2010) The genomic location of sexually antagonistic variation: some cautionary comments. \emph{Evolution} 64:1510--1516.
\item Geeta Arun M, Chechi TS, Meena R, Bhosle SD, Srishti, Prasad NG (2022) Investigating the interaction between interlocus and intralocus sexual conflict using hemiclonal analysis in \emph{Drosophila melanogaster}. \emph{BMC Ecol Evol} 22:38. doi:10.1186/s12862-022-01992-0.
\item Geeta Arun M (2025) Tie that binds -- Hill-Robertson interference can produce signals of sexual antagonism even in the absence of sexually antagonistic selection. \emph{bioRxiv} doi:10.1101/2025.01.08.632022.
\item Geeta Arun M \emph{et al.}\ (2025) \emph{Am Nat} 206(6):552--568. Not accessible to us; title and content not quoted.
\item Havird JC \emph{et al.}\ (2019) Selfish mitonuclear conflict. \emph{Curr Biol} 29:R496--R511.
\item Innocenti P, Morrow EH (2010) The sexually antagonistic genes of \emph{Drosophila melanogaster}. \emph{PLoS Biol} 8:e1000335.
\item Kidwell JF \emph{et al.}\ (1977) Regions of stable equilibria for models of differential selection in the two sexes under random mating. \emph{Genetics} 85:171--183.
\item Kimura M (1962) On the probability of fixation of mutant genes in a population. \emph{Genetics} 47:713--719.
\item Morrow EH, Arnqvist G, Pitnick S (2003) Adaptation versus pleiotropy: why do males harm their mates? \emph{Behav Ecol} 14:802--806.
\item Owen ARG (1953) A genetical system admitting of two distinct stable equilibria under natural selection. \emph{Heredity} 7:97--102.
\item Parsons PA (1961) The initial progress of new genes with viability differences between sexes and with sex linkage. \emph{Heredity} 16:103--107.
\item Rice WR (1984) Sex chromosomes and the evolution of sexual dimorphism. \emph{Evolution} 38:735--742.
\item Rice WR (1987) The accumulation of sexually antagonistic genes as a selective agent promoting the evolution of reduced recombination between primitive sex chromosomes. \emph{Evolution} 41:911--914.
\item Ruzicka F \emph{et al.}\ (2019) bioRxiv preprint 117176; values from the primary text.
\item Samant MA (2022) ``The ties that bind'': investigating the interaction between interlocus and intralocus sexual conflict using \emph{Drosophila melanogaster}. PhD thesis, IISER Mohali.
\item Samant M, Moitra P, Sinha S, Prasad NG. A two-locus haploid model for the resolution of intra-locus sexual conflict (poster, IISER Mohali).
\item Singh, Jain, Geeta Arun \& Prasad (2023) Nature and location of modifier alleles determine the resolution of intralocus sexual conflict. \emph{bioRxiv} doi:10.1101/2023.10.25.564090.
\end{itemize}
